\section{结论}
本文将开放技能包的质量核验拆成可以分别检验的证据层，
并在SkillNet-demo上完成包审查、机械干预、真实模型评价、产物对照和网络回放。
当前结果表明，接口一致、引用可定位、边界处置正确和完整产物合格之间存在实质差别。
对同一批构造任务，较高的逐项检查比例不能替代完整任务成功；
增加技能文本或集合多样性也没有自动形成稳定收益。

实际产出是可回放的核验协议、受控输入与验收器、完整失败记录、来源绑定实现及条件化结果。
它们为可靠决策提供基础，但尚未形成通用技能认证或已证明的网络优化优势。
后续工作应以独立语义参照、真实工具执行和跨来源任务验证连接这些证据层，
用真实错误放行、覆盖、产物收益和成本决定是否采用一种评价或编排方法。
